% Auto-generated by docs/build_integrated_paper.py -- do not edit by hand.
% Standalone scaffolding (preamble, nomenclature, limitations, conclusion,
% \setcounter) stripped for \input into the combined manuscript.
\section{Digital-Twin Fidelity and the RL-Utility Paradox}
\label{sec:results1-digital-twin}

Block~1 establishes the digital-twin side of the central claim --- that predictive fidelity and RL training utility are related but distinct objectives --- by contrasting the compact control-oriented v3 surrogate with the physically informed v3.5 across architecture, calibration, data resolution, rollout accuracy, and runtime; whether the more accurate twin is the better PPO training environment is tested in Results~II.

\paragraph{Roadmap boundary and artifact chain.}
\label{ssec:block1-roadmap-boundary}

Block~1 asks whether the surrogate family is credible as a fast, physically interpretable digital twin. It does not yet ask whether a controller trained on that surrogate transfers to the live BOPTEST RTE. The executed roadmap path is:
\begin{enumerate}
  \item collect and train the v3 direct-\Tsupply{} surrogate using \texttt{evaluation/run\_block1.py collect-data} and \texttt{v3-train};
  \item prepare the 15-minute corpus and run v3.5 Stage~A/B/C inverse calibration using \texttt{prepare-15min} and \texttt{v35-calibrate --preset canonical};
  \item retrain v3 on the same 15-minute corpus using \texttt{v3-train-15min} to bound the calibration claim;
  \item build the Hou-and-Evins reporting tables and speed benchmark using \texttt{build-reports} and \texttt{speed-benchmark}.
\end{enumerate}



\subsection{Direct Supply-Temperature Control Interface}
\label{ssec:block1-direct-tsup}

All Block 1 surrogates share a single \emph{direct supply-temperature} (direct-\Tsupply{}) control parameterization: the policy commands the supply-air temperature setpoint directly, and the surrogate consumes that setpoint as a boundary condition instead of modelling the upstream valve/coil actuator chain. The two-dimensional action \(a=(a_0,a_1)\in[-1,1]^2\) maps affinely to physical commands,
\begin{equation}
  \Tsupply = 18 + \tfrac{a_0+1}{2}\,(35-18)~[^\circ\mathrm{C}],
  \qquad
  u_{\mathrm{fan}} = \tfrac{a_1+1}{2}\in[0,1],
  \label{eq:action-map}
\end{equation}
and every surrogate exposes the same transition signature
\begin{equation}
  f:\ (T_{\mathrm{zone}},T_{\mathrm{amb}},h,d,a_0,a_1)\ \longmapsto\ (\That_{t+1},\widehat{P}_{t+1}),
  \label{eq:tsup-signature}
\end{equation}
with \(\widehat{P}\ge 0\) enforced by a Softplus output and capped at \(P_{\max}=5500~\mathrm{W}\). A single adapter wraps four interchangeable backends behind \eqref{eq:tsup-signature}---legacy v3, raw v3.5, calibrated v3.5, and the hybrid pairing---which is what makes the like-for-like fidelity comparisons in this section well posed: only the internal dynamics change, never the control interface or the observation map.

\begin{table}[pos={!ht}]
\centering
\small
\caption{Modelling assumptions of the direct-\Tsupply{} control interface.}
\label{tab:tsup-assumptions}
\begin{tabularx}{\linewidth}{lX}
\toprule
Assumption & Statement \\
\midrule
Actuator abstraction & Supply-air temperature is commanded directly; valve/coil dynamics are not modelled beyond the Stage A latency term. \\
Command bounds & \(\Tsupply\in[18,35]\,^\circ\)C; fan command \(u_{\mathrm{fan}}\in[0,1]\). \\
Power non-negativity & \(\widehat{P}=\mathrm{Softplus}(\cdot)\ge 0\), capped at \(P_{\max}=5500\)~W. \\
State support & Zone and ambient temperatures normalized to \([-1,1]\) over \([15,35]\) and \([-10,40]\,^\circ\)C. \\
Interface invariance & All four backends obey \eqref{eq:tsup-signature}; only the internal dynamics differ. \\
\bottomrule
\end{tabularx}
\end{table}

\subsection{Control-Oriented Surrogate Architecture (v3)}
\label{ssec:block1-v3}

The v3 surrogate is a compact control-oriented direct-\Tsupply{} model. Its role is not to be the most accurate long-horizon forecaster; its role is to provide smooth and fast rollout dynamics for downstream PPO experiments. In contrast, v3.5 is a physics-informed extension whose Stage B parameter is interpretable as the zone thermal capacitance \(\Czon\); embedding thermodynamic structure directly into the learned model follows the physics-informed reinforcement-learning paradigm~\citep{Hedayat2025PINN}.







The v3 transition model is
\begin{align}
  \Delta \That_{t+1} &= f_{\theta,T}(x_t), &
  \widehat{P}_{t+1} &= f_{\theta,P}(x_t), &
  \That_{t+1} &= T_t+\Delta \That_{t+1}.
\end{align}
It is trained by a multi-horizon supervised objective with a physical-consistency penalty,
\begin{equation}
  \mathcal{L}_{v3}
  = \sum_{k\in\{1,2,4\}} \frac{1}{N}\sum_{i}\big(\That_i(t{+}k)-T_i(t{+}k)\big)^2
  + \lambda_{\mathrm{phys}}\,\mathrm{ReLU}\!\big(|\Delta\That|-\Delta T_{\max}\big),
  \label{eq:v3-loss}
\end{equation}
where the multi-step term (horizons of 1, 2, and 4 transitions) enforces short-horizon rollout consistency and the penalty discourages physically implausible single-step jumps. The update \(\That_{t+1}=T_t+\Delta\That_{t+1}\) carries \emph{no} explicit time step: the learned increment is a per-transition quantity, the design choice revisited in Section~\ref{ssec:block1-stepsize}.



The canonical v3 checkpoint is the early-stopped optimum of a 215-epoch supervised run: the best validation epoch is 185, at which the held-out one-step temperature fit reaches \(R^2=0.991\). This single-step score is deliberately optimistic relative to the multi-step rollout error in Section~\ref{ssec:block1-predictive-diagnostics}, which is the metric that actually governs whether the surrogate can substitute for BOPTEST.



The v3 input vector is intentionally low-dimensional. It keeps only the thermal state, outdoor disturbance, cyclic time embedding, and direct-\Tsupply{} action coordinates required by the control problem. The action coordinate \(a_0\in[-1,1]\) is mapped to an 18--35~$^\circ$C supply-temperature command --- a direct supply-temperature setpoint interface of the kind shown effective for DRL-based room HVAC control~\citep{Sun2024DirectTSup} --- and the power head is constrained through a non-negative output transform. These choices make v3 a control-oriented surrogate rather than a generic black-box forecaster. Its full hyperparameters and layer shapes are listed in Supplementary Table~\suppref{S6}.



\subsection{Physics-Informed Twin and Inverse Calibration (v3.5)}
\label{ssec:block1-v35-calibration}

Unlike the purely black-box v3 head, the v3.5 surrogate is built on a first-principles lumped-capacity zone balance with a single effective parameter, the zone thermal capacitance \(\Czon\),
\begin{equation}
  \Czon\,\frac{\mathrm{d}T_{\mathrm{zone}}}{\mathrm{d}t} = \dot{Q}(x;\phi),
  \qquad \dot{Q}(x;\phi)=q_{\mathrm{scale}}\,q_\phi(x),
  \label{eq:rc-ode}
\end{equation}
where \(\dot{Q}(x;\phi)\) is a learned net heat-flow head (\(q_{\mathrm{scale}}=3000\)~W). HVAC input, envelope coupling, and internal gains are absorbed into \(\dot{Q}\) through the feature vector \(x\); there is no separately parameterized \(1/R\) term. With control step \(\Delta t\), \eqref{eq:rc-ode} is integrated by an explicit Euler step,
\begin{equation}
  \That_{t+1} = \mathrm{clip}\!\Big(T_t + \Delta t\,\frac{\dot{Q}(x;\phi)}{\Czon},\ T_{\min},\,T_{\max}\Big).
  \label{eq:rc-euler}
\end{equation}
To keep the inverse problem well posed, \(\Czon\) is not a free scalar but a positively reparameterized one,
\begin{equation}
  \Czon = c_s\big(\mathrm{softplus}(\rho) + c_{\min}/c_s\big),
  \qquad \Czon \ge c_{\min} = 5\times10^4~\mathrm{J/K},
  \label{eq:czon-reparam}
\end{equation}
with \(\rho\) the learnable log-parameter and \(c_s=10^5\) a fixed scale; this guarantees a physically positive capacitance throughout optimization. Because \(\Czon\) is a property of the building rather than of the data-generating policy, it is expected to remain meaningful across testcases, which is what makes the inverse identification below worthwhile.

The canonical v3.5 is built by two sequential preset runs. \texttt{block1\_15min\_episodeaware} does the temperature-side inverse identification --- Stage~A aligns telemetry and removes latency/bias, Stage~B identifies \(\Czon\) on excitation-rich windows (120 epochs), Stage~C refines residual heads (60 epochs) while preserving the physical parameter. The second preset, \texttt{block1\_15min\_power\_head\_only}, freezes \(\Czon\), skips Stage~B, and refines only the power residual head (80 epochs) --- which is why its JSON reports \texttt{stage\_b\_epochs\_ran=0}.

Stage A estimates a one-step latency compensation, a temperature bias correction of -0.105~$^\circ$C, and a postprocessed temperature RMSE of 0.238~$^\circ$C. Stage B uses the top-excitation subset: 403 of 8058 training rows at excitation quantile 0.95. This filtering is important because \(\Czon\) is identifiable mainly during transients, not quasi-steady operation.





Formally, Stage B is a maximum-a-posteriori estimate of the single physical scalar,
\begin{equation}
  \hat C_{\mathrm{zon}}
  = \arg\min_{C}\ \sum_{i\in\mathcal{E}}\big(\That_i(C)-T_i\big)^2
  + \frac{1}{2\sigma_p^2}\big(C-C^{\mathrm{prior}}\big)^2,
  \label{eq:stageb-map}
\end{equation}
where \(\mathcal{E}\) is the excitation subset and the quadratic term is a weak Gaussian prior centred on \(C^{\mathrm{prior}}\). The excitation filter is justified by identifiability: in quasi-steady operation the sensitivity \(\partial\That/\partial C\to 0\), so steady windows carry almost no information about \(\Czon\) and would merely pull the estimate back to the prior. Restricting \(\mathcal{E}\) to the top-excitation transients (403 of 8058 rows) concentrates the identification where \(\partial\That/\partial C\) is large, which is why the mean excitation score on \(\mathcal{E}\) is several times the corpus mean (Table~\suppref{S9}).

The intermediate Stage~B estimate is \(\Czon=4.413\times10^5~\mathrm{J/K}\), and the relative update over the physical prior is
\begin{equation}
  \eta_C =
  \frac{\Czon^{\mathrm{final}}-\Czon^{\mathrm{prior}}}{\Czon^{\mathrm{prior}}}
  = 0.0506,
\end{equation}
where \(\Czon^{\mathrm{final}}=4.413 \times 10^5~\mathrm{J/K}\). The update remains inside the prior physical plausibility band.

Beyond the point estimate, two complementary uncertainty checks apply. The optimizer is well conditioned (over the final 20 Stage~B epochs the \(\Czon\) trajectory is stable to within \(\pm1,005\)~J/K, 0.23\%). More importantly, a Laplace/Fisher identifiability interval --- obtained from the local curvature of the one-step temperature SSE with respect to \(\Czon\) (residual heads fixed) over the high-excitation transients (\(n=1838\)) --- gives a \(1\sigma\) uncertainty of 6.6\% and a 95\% interval \([3.845,\,4.980]\times10^5\)~J/K. The \(+5.1\%\) data-driven update therefore sits within roughly one standard error of the prior: \(\Czon\) is identifiable with a moderate, honestly bounded signal rather than a tightly pinned value. Physically, the identified capacitance corresponds to an equivalent thermal mass of \(\Czon/c_{p,\mathrm{air}}\approx439\)~kg of air (\(c_{p,\mathrm{air}}=1005~\mathrm{J\,kg^{-1}K^{-1}}\)), or roughly 366~m\(^3\) of equivalent air volume at \(\rho=1.2~\mathrm{kg\,m^{-3}}\). This is the expected order of magnitude for the \texttt{bestest\_air} zone augmented by the effective thermal mass of its internal surfaces, which is exactly why the data-driven update over the air-only prior is small (\(+5\%\)) rather than large.

\medskip\noindent\textbf{Physical-consistency audit.} As a final check we probe the learned dynamics directly: sweeping the supply-temperature command across its range with the state held fixed, the predicted next-step zone temperature increases with \Tsupply{} in 100\% of 400 sampled states (correct sign in every case), with a mean sensitivity of \(+0.061\)~\(^\circ\)C per \(^\circ\)C and a mean Spearman rank correlation of 0.91 (Figure~\suppref{S4}). The response is monotone in the physically expected direction; the small departures from strict monotonicity reflect the neural residual head and stay well inside the surrogate temperature tolerance.




\begin{figure}[pos={!ht}]
  \centering
  \includegraphics[width=\linewidth]{\detokenize{rie_fig03_stage_abc_diagnostics.pdf}}
  \caption{Stage A/B/C calibration diagnostics. Panel (a) shows physical-parameter convergence inside the prior band; panel (b) shows the measured improvement in temperature and power fidelity.}
  \label{fig:stage-abc}

{\footnotesize\itshape Data: \texttt{outputs/\allowbreak surrogate\_\allowbreak v35\_\allowbreak inverse\_\allowbreak boptest\_\allowbreak 15min\_\allowbreak episodeaware/\allowbreak stage\_\allowbreak b\_\allowbreak history\_\allowbreak v35.csv}}
\end{figure}



\begin{table}[pos={!ht}]
\centering
\small
\caption{Key design choices for the v3/v3.5 surrogate family and their engineering rationale. A formal sensitivity sweep over these settings is identified as future work in Section~\ref{ssec:block1-limitations}.}
\label{tab:design-rationale}
\begin{tabularx}{\linewidth}{llX}
\toprule
Choice & Value & Rationale \\
\midrule
Hidden dimension & 64 & smooth, low-curvature response surface for stable policy gradients \\
Optimizer / LR & AdamW / 1e-3 & decoupled weight decay; cosine annealing \\
Batch size & 256 & variance/throughput trade-off on the prepared corpus \\
Heat-flow scale $q_{\mathrm{scale}}$ & $3000$ W & keeps $q_\phi(x)$ near unit range for conditioning \\
Capacitance floor $c_{\min}$ & $5\times10^4$ J/K & physical lower bound enforced by reparameterization \\
Excitation quantile & 0.95 & isolates transient windows where $\partial \widehat{T}/\partial C$ is large \\
$\Czon$ scalar LR & $10^{-3}$ & dedicated rate for the single physical parameter \\
\bottomrule
\end{tabularx}
\end{table}

\subsection{Multi-Step Validation and Matched-Corpus Ablation}
\label{ssec:block1-predictive-diagnostics}

\subsubsection{Multi-step predictive validation}

For horizon \(h\), the rollout error and RMSE are
\begin{align}
  e_i(h) &= \That_i(t+h)-T_i(t+h),\\
  \RMSE_T(h) &= \sqrt{\frac{1}{N}\sum_{i=1}^N e_i(h)^2}.
\end{align}
Engineering tolerance is reported as an absolute-error CDF,
\begin{equation}
  F_{|e|}(\epsilon)=\Pr(|\That-T_{\mathrm{BOPTEST}}|\leq \epsilon).
\end{equation}
The coefficient of determination follows the standard definition \(R^2 = 1 - \mathrm{SS}_{\mathrm{res}}/\mathrm{SS}_{\mathrm{tot}}\) against the held-out mean, so a \emph{negative} 24 h \(R^2\) (as for the hourly v3 in Table~\suppref{S11}) means the long-horizon rollout is worse than predicting that mean. The error tail is summarized by the 95th-percentile absolute error \(\mathrm{P95}=\inf\{\epsilon:F_{|e|}(\epsilon)\ge 0.95\}\).



The calibrated v3.5 surrogate reaches a 24 h rollout \(\RMSE_T=0.644\)~$^\circ$C, compared with raw v3.5 at 1.466~$^\circ$C and the canonical v3 hourly reference at 1.557~$^\circ$C.



Table~\suppref{S11} exposes the split between the two families. The hourly v3 has a \emph{negative} 24~h \(R^2\) --- worse than the historical mean over long horizons, and worse than a naive persistence forecast at 24~h (1.557 vs 0.969~$^\circ$C) --- yet it remains a usable control surrogate because PPO needs smooth local gradients, not long-horizon realism. Calibrated v3.5 instead keeps a positive 24~h \(R^2\), bounds its tail (24~h P95 1.207~$^\circ$C), and beats persistence at every horizon, confirming v3 is a control surrogate rather than a forecaster.



The episode-level diagnostic is important because a low aggregate RMSE can hide a single unstable rollout. The calibrated v3.5 artifact remains below the 1~$^\circ$C reference in all eight held-out episodes; its best episode is peak heat thermostatic at 0.486~$^\circ$C and its worst episode is typical heat hdrl at 0.964~$^\circ$C. This supports predictive validity as a repeated-episode result rather than a single-window artifact.

\subsubsection{Matched-corpus ablation}
\label{ssec:block1-corpus-matched}

The roadmap includes a corpus-matched v3 retraining step because the original v3 and canonical v3.5 artifacts use different data resolutions. This is not a cosmetic check: without it, a reviewer could correctly attribute the v3.5 gain to the 15-minute corpus rather than Stage A/B/C calibration.

\begin{table}[pos={!ht}]
\centering
\small
\caption{Corpus-matched comparison.}
\label{tab:matched}
\begin{tabular}{lrlr}
\toprule
Variant & Step (s) & Calibration & 24 h \(\RMSE_T\) \\
\midrule
v3 hourly & 3600 & n/a (control-oriented black box) & 1.557 \\
v3 15-min matched & 900 & n/a (control-oriented black box) & 0.876 \\
raw v3.5 & 900 & OFF (no Stage A/B/C) & 1.466 \\
calibrated v3.5 & 900 & FULL (Stage A + B + C, C\_zon = 4.413e5 J/K) & 0.644 \\
\bottomrule
\end{tabular}
\end{table}

\begin{figure}[pos={!ht}]
  \centering
  \includegraphics[width=\linewidth]{\detokenize{rie_fig05_matched_corpus_attribution.pdf}}
  \caption{Matched-corpus attribution. The v3-to-v3.5 gain is split into the 15-minute corpus contribution and the Stage A/B/C calibration contribution.}
  \label{fig:matched}

{\footnotesize\itshape Data: \texttt{reports/\allowbreak block1\_\allowbreak corpus\_\allowbreak matched\_\allowbreak comparison.csv}}
\end{figure}

The attribution is an additive split of the legacy-v3-to-calibrated-v3.5 gap,
\begin{equation}
  \underbrace{\Delta_{\mathrm{total}}}_{\text{v3 hourly}\,\to\,\text{cal. v3.5}}
  = \underbrace{\Delta_{\mathrm{corpus}}}_{\text{v3 hourly}\,\to\,\text{v3 15-min}}
  + \underbrace{\Delta_{\mathrm{cal}}}_{\text{v3 15-min}\,\to\,\text{cal. v3.5}},
  \label{eq:decomp}
\end{equation}
which holds because the matched-architecture path is additive (the intermediate point is the same v3 retrained on the 15-min corpus). Of the total 0.913~$^\circ$C drop from legacy v3 to calibrated v3.5, the matched-corpus report attributes 0.681~$^\circ$C (74.6\%) to the corpus shift and 0.232~$^\circ$C (25.4\%) to Stage~A/B/C. (An alternative split through the raw v3.5 backbone gives 9.9\% corpus / 90.1\% calibration; the two differ because corpus and architecture interact non-additively, and we quote the matched-architecture split since both endpoints are control-architecture surrogates.) Either way the calibration claim is positive but bounded.

\subsection{Runtime Feasibility}
\label{ssec:block1-runtime}

The final Block 1 roadmap command builds the speed benchmark. The benchmark compares the same BOPTEST RTE HTTP loop used by the paper against local surrogate stepping under a 900 s control step. The result is not a controller claim; it only establishes feasibility for downstream policy optimization.





At the hybrid throughput, a single \(5\times10^6\)-step policy-optimization run costs about 46.6 minutes of pure environment time, versus roughly 66.1 hours on the live BOPTEST RTE HTTP loop. The calibrated v3.5 surrogate is 114.2x faster than the BOPTEST RTE HTTP loop, and the hybrid reference remains 85.0x faster. This closes the runtime part of Block~1: the surrogate family is not only more accurate after calibration, but also fast enough to support repeated policy-training experiments without using the RTE HTTP loop as the inner training environment.

\subsection{Step-Size Design Disclosure}
\label{ssec:block1-stepsize}

The canonical v3 checkpoint is trained at a 3600~s step but applied at the 900~s control step; since its transition \(\That_{t+1}=T_t+\Delta\That_{t+1}\) carries no explicit timestep scaling, the hourly increment is used unchanged --- the structural reason the legacy v3 carries a 1.557~$^\circ$C 24~h error while the corpus-matched v3 reaches 0.876~$^\circ$C (Table~\ref{tab:matched}). The mismatch is preserved deliberately: every controller-facing KPI is measured on the live BOPTEST simulator (the surrogate only supplies a smooth gradient signal), and swapping checkpoints would break the frozen audit chain of the canonical experiments. The matched-corpus model grounds the attribution here and is evaluated downstream as a closed-loop ablation in Results~II (Section~\ref{sec:results2-control}), testing whether v3's training utility comes from temporal coarse-graining rather than its black-box architecture. Unlike v3, the v3.5 forward pass \eqref{eq:rc-euler} carries an explicit \(\Delta t\), so a warm-started v3.5 rescales the inherited v3 heat term by
\begin{equation}
  s = s_{\mathrm{legacy}}\,\frac{C_{\mathrm{init}}}{\Delta t_{\mathrm{legacy}}\,q_{\mathrm{scale}}},
  \label{eq:warmstart-scale}
\end{equation}
so that the per-step increment is dimensionally consistent at the 900~s runtime step. This time-scaling is exactly what the v3 update lacks.
